\begin{Exercise}[label={ex:cinema}]

(TdE Luglio 2018) Si considerino le seguenti frasi, esprimibili in logica proposizionale tramite tre predicati (A, B, C), che esprimono due condizioni per cui si pu\`o entrare al cinema gratis:

\begin{enumerate}
	\item Sei il millesimo cliente E sei il cliente pi\`u paffutello della sala, OPPURE non hai mai visto il film in programma.
	\item Hai gi\`a visto il film in programma E almeno una delle seguenti condizioni \`e vera:
\begin{itemize}
	\item non sei il millesimo cliente;
	\item c'\`e qualche cliente in sala paffutello almeno quanto te.
\end{itemize}
\end{enumerate}

Si definiscano i tre predicati logici A, B, C contenuti nelle due precedenti frasi, si traducano quindi le due frasi in espressioni booleane, si dica se le proposizioni espresse nelle due frasi sono equivalenti, giustificando la risposta.		
		
\end{Exercise}

\begin{Answer}[ref={ex:cinema}]

	Definiamo:
	\begin{itemize}
		\item A:  ``sei il millesimo cliente'';
		\item B: ``sei il cliente pi\`u paffutello della sala'';
		\item C: ``hai visto il film in programma''.
	\end{itemize}
	
	La prima frase corrisponde a: (A \&\& B) || !C.
	
	La seconda frase corrisponde a: C \&\& (!A || !B).
	
	Le due espressioni non sono equivalenti perch\'e, quando C \'e falsa, la prima \'e sempre vera, mentre la seconda \'e sempre falsa. In alternativa si scrivano le tabelle di verit\`a delle due espressioni e si noti che non sono identiche.

\end{Answer}